\documentclass[10pt,a4paper]{amsart}
\usepackage[margin=19mm]{geometry}
\usepackage{amsmath,amssymb,mathtools}
\usepackage[scr=rsfs,cal=euler]{mathalfa}
\usepackage{xfrac}
\usepackage[unicode,hidelinks,bookmarks=false]{hyperref}
\newcommand{\ie}{i.e.\ }
\newcommand{\cf}{cf.\ }
\newcommand{\resp}{respectively\ }
\newcommand{\Subsection}{\subsection}
\providecommand{\nobreakdash}{\nobreak\hbox{-}}
\setlength{\emergencystretch}{2em}
\multlinegap0pt
\usepackage{amssymb}
\usepackage{mathtools}
\usepackage{enumerate}
\usepackage{tikz}
\newtheorem{prop}{Proposition}
\newcommand{\R}{\mathbb R}
\title{Runge type approximation results for spaces of smooth Whitney jets}
\author{Tomasz Cia\'s, Thomas Kalmes}
\date{Source: arXiv:2502.03815v3 (CC BY 4.0)}
\begin{document}
\maketitle

\noindent\emph{Source and licence.} Runge type approximation results for spaces of smooth Whitney jets. Version arXiv:2502.03815v3 (CC BY 4.0), distributed by arXiv under the Creative Commons Attribution 4.0 International licence, http://creativecommons.org/licenses/by/4.0/, as recorded in the arXiv licence metadata for this exact version and shown on the abstract page https://arxiv.org/abs/2502.03815v3. Full licence notice: \texttt{licenses/arxiv-2502.03815-CC-BY-4.0.txt}.\par\smallskip

\noindent\textit{Editorial scope.} Counted unit: the article's prop d>2 (source0.tex lines 655--659). The article's complete original proof of it is delivered with it (source0.tex lines 661--740, one \texttt{proof} environment of 80 lines). No mathematics is written, repaired or re-derived: the statement and the proof are the article's own lines, in the article's own notation, and no retained line is substituted or re-flowed.  No further source-local context is needed: every cross-reference of every retained line resolves inside the two delivered spans.  Nothing was omitted from any retained span, so the package carries no omission record.  No retained line contains a citation, so the wrapper carries no bibliography and no bibliography entry.  No retained line contains a figure environment, an image-inclusion command or drawing code, so no image is delivered and the package declares no asset dependency.  Editorial formatting only, in addition: an AMS \texttt{amsart} preamble that restates the article's own declarations and macros used by the retained lines, namely the environments \texttt{prop} on separate counters, together with the standard packages those lines need.  The retained environments print in this document as Proposition 1 in order of appearance, whereas the article prints the same items with the numbering of its own full document.  The complete native source of the article as distributed in the arXiv e-print for this exact version is bundled unchanged as \texttt{sources/172831/source0.tex}.
\begin{prop}\label{prop:C1-boundary, d>2}
	Assume that $d\geq 3$ and that $F\subset\R^d$ is closed such that $\partial F$ is a $C^1$-hypersurface with the property that for every $\xi\in \partial F$ its normal space is not spanned by $e_1$. 
	
	Moreover, let $c\in\R$ be such that $\left(\R^d\backslash F\right)\cap \{c\}\times\R^{d-1}$ has a bounded connected component $C$. Then, there is $\delta>0$ such that the connected component of $\left(\R^d\backslash F\right)\cap (c-\delta,c+\delta)\times\R^{d-1}$ which contains $C$, is bounded, too. 
\end{prop}

\begin{proof}
Let $c\in\R$ and $C\subset \left(\R^d\backslash F\right)\cap \{c\}\times\R^{d-1}$ be as in the hypothesis. We denote the boundary of $C$ in $\{c\}\times\R^{d-1}$ by $\partial_c C$. Then, $\partial_c C$ is a compact subset of $\{c\}\times\R^{d-1}$ and $\left(\{c\}\times\R^{d-1}\right)\setminus \partial_c C$ has exactly one unbounded connected component which we denote by $D_\infty$. In particular, $D_\infty$ is open in $\{c\}\times\R^{d-1}$, disjoint to $C$, and 
\begin{equation}\label{eq:set inclusion 1}
	\partial_cD_\infty\subset\partial_cC\subset\left(\{c\}\times\R^{d-1}\right)\cap \partial F,
\end{equation}
where $\partial_c D_\infty$ is the boundary of $D_\infty$ in $\{c\}\times\R^{d-1}$. 

Let us fix $\xi\in\partial_c D_\infty$. By hypothesis, there is an open, connected subset $U_\xi$ of $\R^d$ conataining $\xi$ and a $C^1$-function $f_\xi:U_\xi\rightarrow \R$ such that
	$$\partial F\cap U_\xi=\{x\in U_\xi \colon f_\xi(x)=0\}$$
	and $e_1, \nabla f_\xi(x)$ are linearly independent for each $x\in U_\xi$. 
Then
$\frac{\partial f_\xi}{\partial x_j}(\xi)\neq0$ for some $2\leq j\leq d$, and w.l.o.g. we may assume that $\frac{\partial f_\xi}{\partial x_d}(\xi)\neq0$. 
From the implicit function theorem it follows that there is $\delta_\xi>0$ for which the open sets
$$Q_\xi:=(\xi_1-\delta_\xi,\xi_1+\delta_\xi)\times\ldots\times(\xi_{d-1}-\delta_\xi,\xi_{d-1}+\delta_\xi)$$
and
$$P_\xi:=(\xi_d-\delta_\xi,\xi_d+\delta_\xi)$$
satisfy $V_\xi:=Q_\xi\times P_\xi\subset U_\xi$ and there is a unique
$C^1$-function $g_\xi\colon Q_\xi\to P_\xi$,
such that
$$g_\xi(\xi_1,\ldots,\xi_{d-1})=\xi_d\quad\text{and}\quad f_\xi(x,g_\xi(x))=0$$
for all $x\in Q_\xi$.
Then 
$$\partial F\cap V_\xi=\{(x,y)\in Q_\xi\times P_\xi\colon f_\xi(x,y)=0\}=\{(x,g_\xi(x))\colon x\in Q_\xi\},$$ 
so
$$V_\xi\setminus\partial F=V_\xi^-\cup V_\xi^+,$$
where the sets
$$V_\xi^-:=\{(x,y)\in Q_\xi\times P_\xi\colon y<g_\xi(x)\}\quad\text{and}\quad
V_\xi^+:=\{(x,y)\in Q_\xi\times P_\xi\colon y>g_\xi(x)\}$$
are non-empty, pairwise disjoint, open and connected. Therefore, the set $V_\xi\setminus\partial F$ has exactly two open connected components $V_\xi^-$ and $V_\xi^+$.

Similarly, the set $(V_\xi\cap\{c\}\times\R^{d-1})\setminus\partial F$ is divided by the graph of a $C^1$-function $h_\xi\colon \widehat{Q_\xi}\to P_\xi$, $h_\xi(x_2,\ldots,x_{d-1})=g_\xi(\xi_1,x_2,\ldots,x_{d-1})$ into two non-empty, open, connected sets, where
$$\widehat{Q_\xi}=(\xi_2-\delta_\xi,\xi_2+\delta_\xi)\times\ldots\times(\xi_{d-1}-\delta_\xi,\xi_{d-1}+\delta_\xi),$$
and the decomposition is the following
$$(V_\xi\cap\{c\}\times\R^{d-1})\setminus\partial F=
(V_\xi^-\cap\{c\}\times\R^{d-1})\cup (V_\xi^+\cap\{c\}\times\R^{d-1}).$$

Let us note that the sets $V_\xi\cap D_\infty$ and $V_\xi\cap C$ are pairwise disjoint non-empty open subsets of $V_\xi\cap\{c\}\times\R^{d-1}$. 
We claim that $(V_\xi\cap\{c\}\times\R^{d-1})\setminus\partial F\subset C\cup D_\infty$. Indeed, if this inclusion were not true, then we would find a bounded connected component $B$ of $\left(\{c\}\times\R^{d-1}\right)\setminus \partial_c C$ which is different from $C$ and such that 
$V_\xi\cap B$ is a non-empty open set. Then $(V_\xi\cap\{c\}\times\R^{d-1})\setminus\partial F$ would have at least three connected components as none two sets among $V_\xi\cap D_\infty$, $V_\xi\cap C$ and $V_\xi\cap B$ are in the same connected component, a contradiction. The inclusion just demonstrated leads easily to the identity  
$$(V_\xi\cap\{c\}\times\R^{d-1})\setminus\partial F=(V_\xi\cap C)\cup (V_\xi\cap D_\infty).$$
Since, $(V_\xi\cap\{c\}\times\R^{d-1})\setminus\partial F$ has exactly two open connected components, we have
$$V_\xi\cap C=V_\xi^-\cap\{c\}\times\R^{d-1}\quad\text{and}\quad V_\xi\cap D_\infty=V_\xi^+\cap\{c\}\times\R^{d-1}$$ 
or 
$$V_\xi\cap C=V_\xi^+\cap\{c\}\times\R^{d-1}\quad\text{and}\quad V_\xi\cap D_\infty=V_\xi^-\cap\{c\}\times\R^{d-1}.$$ 

Since $\partial_c D_\infty\subset \partial_c C$ we have that $\partial_c D_\infty$ is compact and $\{V_\xi:\xi\in\partial_c D_\infty\}$ is an open covering of $\partial_c D_\infty$. Additionally, for every $\xi\in \partial_c D_\infty$ it holds
	$$V_\xi\cap\partial F\cap\left(\{c\}\times\R^{d-1}\right)=V_\xi\cap \partial_c D_\infty.$$
Indeed, ''$\supset$'' follows directly from \eqref{eq:set inclusion 1}. In order to show ''$\subset$'', we observe
\begin{eqnarray*}
	\overline{V_\xi^+\cap(\{c\}\times\R^{d-1})}^{\{c\}\times\R^{d-1}}&=&\overline{\{(c,x_2,\ldots,x_d): x_d> g_\xi(c,x_2,\ldots,x_{d-1})\}}^{\{c\}\times\R^{d-1}}\\
	&\supset& \{(c,x_2,\ldots,x_d): x_d= g_\xi(c,x_2,\ldots,x_{d-1})\}\\
	&=&V_\xi\cap \partial F\cap(\{c\}\times\R^{d-1})
\end{eqnarray*}
where closures are taken in $\{c\}\times\R^{d-1}$. Thus, in case of $V_\xi\cap D_\infty=V_\xi^+\cap(\{c\}\times\R^{d-1})$ we conclude
$$V_\xi\cap \partial F\cap(\{c\}\times\R^{d-1})\subset \overline{V_\xi\cap D_\infty}^{\{c\}\times\R^{d-1}}\subset \overline{D_\infty}^{\{c\}\times\R^{d-1}}=\partial_c D_\infty\cup D_\infty$$
which by $\partial F\cap D_\infty=\emptyset$ implies
$$V_\xi\cap\partial F\cap(\{c\}\times\R^{d-1})\subset V_\xi\cap \partial_c D_\infty.$$
Analogously, in case of $V_\xi\cap D_\infty=V_\xi^-\cap(\{c\}\times\R^{d-1})$, one uses with the same arguments
$$\overline{V_\xi^-\cap(\{c\}\times\R^{d-1})}^{\{c\}\times\R^{d-1}}\supset V_\xi\cap \partial F\cap(\{c\}\times\R^{d-1})$$
to conclude again $V_\xi\cap\partial F\cap(\{c\}\times\R^{d-1})\subset V_\xi\cap \partial_c D_\infty$, as desired.

Let $\xi^{(1)},\ldots,\xi^{(m)}\in\partial_c D_\infty$ be such that $\{V_{\xi^{(1)}},\ldots, V_{\xi^{(m)}}\}$ is an open covering of $\partial_c D_\infty$. We define
	$$\delta:=\min\{\delta_{\xi^{(j)}}: 1\leq j\leq  m\}\text{ and }  V_j:=V_{\xi^{(j)}}\cap((c-\delta,c+\delta)\times\R^{d-1})$$
for $j=1,\ldots,m$.
Then, $\{V_1,\ldots, V_m\}$  
remains an open covering of $\partial _c D_\infty$ and $B:=\cup_{j=1}^m V_j\cap \partial F$ is a subset of $(c-\delta,c+\delta)\times \R^{d-1}$ which satisfies
$$B\cap \{c\}\times\R^{d-1}=\cup_{j=1}^m V_j\cap \partial F\cap(\{c\}\times\R^{d-1})=\cup_{j=1}^m V_j\cap \partial_c D_\infty=\partial_c D_\infty.$$
Let $C^\delta$ be the connected component of $(\R^d\backslash B)\cap (c-\delta,c+\delta)\times\R^{d-1}$ which contains $C$. Additionally, since $D_\infty$ is connected and obviously contained in any unbounded connected component of $(\R^d\backslash B)\cap (c-\delta,c+\delta)\times\R^{d-1}$, 
there is precisely one such component. We denote it by $D_\infty^\delta$. Then, denoting the boundary of $D_\infty^\delta$ in $(c-\delta,c+\delta)\times \R^{d-1}$ by $\partial_\delta D_\infty^\delta$ it holds
	$$B\cap\{c\}\times\R^{d-1}=\partial_c D_\infty\subset \partial_\delta D_\infty^\delta\subset\partial F.$$ 

Without loss of generality we may assume that $\delta_{\xi_1}=\delta$ which gives $V_1=V_{\xi_1}$, and also that $V_1\cap C=V_1^-\cap\{c\}\times\R^{d-1}$ and $V_1\cap D_\infty=V_1^+\cap\{c\}\times\R^{d-1}$. Then $V_1^-$ is a connected subset of 
$$(\R^d\backslash B)\cap (c-\delta,c+\delta)\times\R^{d-1}$$ which intersects $C$. Therefore, $V_1^-\subset C^\delta$ so that $V_1^-\subset V_1\cap C^\delta$. Similarly, $V_1^+\subset V_1\cap D_\infty^\delta$. From $V_1\setminus\partial F=V_1^+\cup V_1^-$ and
$$D_\infty^\delta\cup C^\delta\subset (\R^d\backslash B)\cap (c-\delta,c+\delta)\times\R^{d-1}$$
we derive $V_1\cap C^\delta\subset V_1^+\cup V_1^-$ and $V_1\cap D_\infty^\delta\subset V_1^+\cup V_1^-$. Combining these inclusions with the disjointness of $C^\delta$ and $D_\infty^\delta$ as well as  $V_1^-\subset V_1\cap C^\delta$ and $V_1^+\subset V_1\cap D_\infty^\delta$, respectively, yields $V_1\cap C^\delta=V_1^-$ and $V_1\cap D_\infty^\delta=V_1^+$, respectively.

In order to complete the proof, we have to show that $C^\delta$ and $D_\infty^\delta$ are disjoint. But this is not the case, as $C^\delta=D_\infty^\delta$ gives
$$V_1^-=V_1\cap C^\delta=V_1\cap D_\infty^\delta=V_1^+,$$
a contradiction. 
\end{proof}

\end{document}
